% Source: https://artofproblemsolving.com/community/c3h1216440_1st_reun_mock_middle_school_math_contest
% Stable identity: aops-topic-1216440-q18
Both bases have radius $3$. The hemisphere's total surface area is its curved part plus its flat disc:

\begin{asy}
usepackage("amsmath");
size(6cm);
defaultpen(fontsize(10pt));
draw((0,0)--(1.5,2.598)--(3,0)); draw(ellipse((1.5,0),1.5,0.35)); draw((1.5,0)--(1.5,2.598),dashed); label("$h$",(1.5,1.3),E); label("$3$",(2.25,0),N);
draw(arc((6,0),1.5,0,180)); draw(ellipse((6,0),1.5,0.35)); draw((6,0)--(7.5,0)); label("$3$",(6.75,0),N);
label("cone",(1.5,-0.5),S); label("hemisphere",(6,-0.5),S);
label("\mbox{}",(min(currentpicture,user=true)+max(currentpicture,user=true))/2,black); // keeps the text after the figure black when LaTeX typesets the labels inline
\end{asy}

\[2\pi\cdot3^2+\pi\cdot3^2=27\pi.\]
The cone's total surface area is $\pi r^2+\pi r\ell=9\pi+3\pi\ell$, where $\ell$ is the slant height. Setting $9\pi+3\pi\ell=27\pi$ gives $\ell=6$, and then the height is
\[h=\sqrt{\ell^2-r^2}=\sqrt{36-9}=\boxed{3\sqrt3}\ \text{inches}.\]
